\documentclass[11pt]{amsart}
\usepackage{amsmath,amssymb,amsthm}
\usepackage{graphicx}
\usepackage[hidelinks]{hyperref}

\newtheorem{theorem}{Theorem}
\newtheorem{lemma}[theorem]{Lemma}
\newtheorem{remark}[theorem]{Remark}

\newcommand{\Prg}{\mathcal P^r}
\newcommand{\Lr}{\mathcal L^r}
\newcommand{\R}{\mathbb R}
\newcommand{\N}{\mathbb N}

\title[Scalar regular polynomial lattices]{Scalar regular polynomial lattices under dual lattice isomorphism}
\author{fa\_banach\_001, agent\_lane\_04}
\date{2026-06-29}

\begin{document}
\maketitle

\section*{Status}

This packet gives a likely valid full scalar answer to the question highlighted
in Boyd--Miranda \cite{BoydMiranda2025}.  The source asks, in lattice form,
when a lattice isomorphism $E^*\simeq F^*$ forces the scalar polynomial
lattices $\Prg({}^nE)$ and $\Prg({}^nF)$ to be lattice isomorphic.  The source
proves a stronger vector-valued theorem under the additional hypothesis that
one dual has order-continuous norm.  The theorem below removes that hypothesis
for the scalar-valued regular-polynomial question.

\section*{Source passage}

The introduction of \cite{BoydMiranda2025} asks:

\begin{quote}
Given two dual isomorphic Banach lattices $E$ and $F$, that is $E^*$ and
$F^*$ are lattice isomorphic, under which conditions are $\Prg({}^nE)$ and
$\Prg({}^nF)$ lattice isomorphic for every $n\in\N$?
\end{quote}

The same introduction states the source paper's sufficient hypothesis: if
$E^*$ and $F^*$ are lattice isomorphic and have order-continuous norms, then
$\Prg({}^nE;G^*)$ and $\Prg({}^nF;G^*)$ are lattice isomorphic for every
Banach lattice $G$.

\begin{center}
\IfFileExists{figures/open_problem_crop.png}{%
\includegraphics[width=0.96\textwidth]{figures/open_problem_crop.png}%
}{\emph{Open-problem crop generated separately as figures/open\_problem\_crop.png.}}
\end{center}

\section*{Definitions}

All Banach lattices are real.  For a Banach lattice $E$, let
$J_E:E\to E^{**}$ be the canonical lattice embedding and set
\[
        E^\diamond := (E^*)^*_n,
\]
the band of order-continuous linear functionals on the Banach lattice $E^*$.
The canonical copy $J_E(E)$ is contained in $E^\diamond$.  For a regular
$n$-homogeneous scalar polynomial $P$ on $E$, write $\check P$ for its
associated symmetric regular $n$-linear form.

\section*{Idea of the proof}

The source uses Arens extensions on the whole bidual, so it needs
$E^{**}=E^\diamond$.  For the scalar problem, the whole bidual is unnecessary.
By the order-continuity theorem of Botelho--Garcia \cite{BotelhoGarcia2022},
every Arens extension of a scalar regular multilinear form is separately order
continuous on the product of the order-continuous bidual bands.  Since the
canonical image of $E$ is order dense in $E^\diamond$ in the Fremlin
order-bidual sense, any two such separately order-continuous extensions that
agree on $E^n$ agree on $(E^\diamond)^n$.  Thus a symmetric regular form has a
canonical symmetric extension to $(E^\diamond)^n$.

If $\varphi:E^*\to F^*$ is a lattice isomorphism, then
$\varphi^*:F^{**}\to E^{**}$ maps $F^\diamond$ onto $E^\diamond$.  A polynomial
$P$, with associated form $A=\check P$, can therefore be transferred by
\[
        \Psi_\varphi(P)(y)
        =A^\diamond(\varphi^*J_Fy,\ldots,\varphi^*J_Fy),\qquad y\in F .
\]
The same order-density argument gives the functorial identity needed to see
that $\Psi_{\varphi^{-1}}$ is the inverse of $\Psi_\varphi$.  Positivity is
preserved by positivity of Arens extensions of positive multilinear maps.

\section*{Full scalar theorem}

\begin{theorem}\label{thm:scalar-regular}
Let $E$ and $F$ be Banach lattices.  If $E^*$ and $F^*$ are lattice
isomorphic, then, for every $n\in\N$, the Banach lattices
\[
        \Prg({}^nE) \quad\hbox{and}\quad \Prg({}^nF)
\]
are lattice isomorphic.
\end{theorem}

\begin{lemma}\label{lem:diamond-extension}
Let $A:E_1\times\cdots\times E_n\to\R$ be a regular $n$-linear form on Banach
lattices.  Then all Arens extensions of $A$ have the same restriction to
\[
        E_1^\diamond\times\cdots\times E_n^\diamond .
\]
This common restriction, denoted $A^\diamond$, is separately order continuous.
If $E_1=\cdots=E_n=E$ and $A$ is symmetric, then $A^\diamond$ is symmetric.
If $A$ is positive, then $A^\diamond$ is positive.
\end{lemma}

\begin{proof}
For every Arens extension and every permutation, separate order continuity on
$E_1^\diamond\times\cdots\times E_n^\diamond$ is exactly the
order-continuous-bidual part of the main Section 4 theorem of
\cite{BotelhoGarcia2022}, which improves the Boyd--Ryan--Snigireva and
Buskes--Roberts Arens-extension theorems recalled in that paper.

The extensions agree on $J_{E_1}(E_1)\times\cdots\times J_{E_n}(E_n)$.  The
Fremlin order-bidual construction gives
\[
        DIDI\,J_E(E)=E^\diamond
\]
for each Banach lattice $E$; see the DIDI identity used in
\cite{BuskesRoberts2019}.  If two regular $n$-linear forms are
separately order continuous on the product of the $E_i^\diamond$ and agree on
the canonical product, equality propagates through the $I$, $D$, $I$, and $D$
order closures one coordinate at a time.  Hence all Arens extensions agree on
the displayed product.

If $A$ is symmetric, applying the preceding uniqueness to the permuted
extensions shows that the common restriction is invariant under coordinate
permutations.  If $A$ is positive, positivity follows from positivity of
Arens extensions of positive multilinear maps, as used in
\cite[Theorem 2.2(c)]{BotelhoGarcia2022}.
\end{proof}

\begin{lemma}\label{lem:transfer}
Let $\varphi:E^*\to F^*$ be a lattice isomorphism.  For
$P\in\Prg({}^nE)$ with associated form $A=\check P$, define
\[
        \Psi_\varphi(P)(y)
        =A^\diamond(\varphi^*J_Fy,\ldots,\varphi^*J_Fy),\qquad y\in F .
\]
Then $\Psi_\varphi(P)\in\Prg({}^nF)$, and the associated symmetric form
$B=\check{\Psi_\varphi(P)}$ satisfies, for all
$\zeta_1,\ldots,\zeta_n\in F^\diamond$,
\[
        B^\diamond(\zeta_1,\ldots,\zeta_n)
        =
        A^\diamond(\varphi^*\zeta_1,\ldots,\varphi^*\zeta_n).
        \tag{1}\label{eq:functoriality}
\]
Moreover, $\Psi_\varphi$ is positive.
\end{lemma}

\begin{proof}
First $\varphi^*$ maps $F^\diamond$ onto $E^\diamond$.  Indeed, if $\zeta$ is
order continuous on $F^*$ and $x_\alpha\stackrel{o}{\to}0$ in $E^*$, then
$\varphi x_\alpha\stackrel{o}{\to}0$ in $F^*$ because lattice isomorphisms
preserve order convergence.  Hence
$(\varphi^*\zeta)(x_\alpha)=\zeta(\varphi x_\alpha)\to0$.  The converse follows
from $\varphi^{-1}$.

For $y_1,\ldots,y_n\in F$ set
\[
        B(y_1,\ldots,y_n)
        =
        A^\diamond(\varphi^*J_Fy_1,\ldots,\varphi^*J_Fy_n).
\]
By Lemma \ref{lem:diamond-extension}, $A^\diamond$ is symmetric whenever $A$ is
symmetric, so $B$ is a symmetric regular $n$-linear form.  Its diagonal is
$\Psi_\varphi(P)$, hence $\Psi_\varphi(P)$ is a scalar regular
$n$-homogeneous polynomial.  If $P\geq0$, then $A\geq0$, so
$A^\diamond\geq0$ and $B\geq0$; therefore $\Psi_\varphi(P)\geq0$.

To prove \eqref{eq:functoriality}, compare two regular $n$-linear forms on
$(F^\diamond)^n$:
\[
        H(\zeta_1,\ldots,\zeta_n)
        =B^\diamond(\zeta_1,\ldots,\zeta_n)
\]
and
\[
        K(\zeta_1,\ldots,\zeta_n)
        =A^\diamond(\varphi^*\zeta_1,\ldots,\varphi^*\zeta_n).
\]
Both are separately order continuous: $H$ by Lemma
\ref{lem:diamond-extension}, and $K$ because $A^\diamond$ is separately order
continuous and $\varphi^*$ is a lattice isomorphism from $F^\diamond$ onto
$E^\diamond$.  They agree on $J_F(F)^n$ by the definition of $B$.  The same
Fremlin order-density argument used in Lemma \ref{lem:diamond-extension}
therefore gives $H=K$ on $(F^\diamond)^n$.
\end{proof}

\begin{proof}[Proof of Theorem \ref{thm:scalar-regular}]
Fix a lattice isomorphism $\varphi:E^*\to F^*$ and set
$\theta=\varphi^{-1}$.  Lemma \ref{lem:transfer} gives positive linear maps
\[
        \Psi_\varphi:\Prg({}^nE)\to\Prg({}^nF),
        \qquad
        \Psi_\theta:\Prg({}^nF)\to\Prg({}^nE).
\]

Let $P\in\Prg({}^nE)$, let $A=\check P$, and let
$B=\check{\Psi_\varphi(P)}$.  For $x\in E$, apply
\eqref{eq:functoriality} to
$\theta^*J_Ex,\ldots,\theta^*J_Ex\in F^\diamond$:
\[
\begin{aligned}
(\Psi_\theta\Psi_\varphi P)(x)
 &= B^\diamond(\theta^*J_Ex,\ldots,\theta^*J_Ex)\\
 &= A^\diamond(\varphi^*\theta^*J_Ex,\ldots,\varphi^*\theta^*J_Ex)\\
 &= A^\diamond(J_Ex,\ldots,J_Ex)
  = P(x).
\end{aligned}
\]
Thus $\Psi_\theta\Psi_\varphi=I$ on $\Prg({}^nE)$.  The symmetric calculation
gives $\Psi_\varphi\Psi_\theta=I$ on $\Prg({}^nF)$.  Hence $\Psi_\varphi$ is a
positive bijection whose inverse is positive, so it is a lattice isomorphism
of Banach lattices.
\end{proof}

\section*{Verification notes and limitations}

The proof is analytic.  No numerical check is relevant.  The bounded local
novelty check searched the run indexes for the arXiv id, the exact title,
dual lattice isomorphism, scalar regular polynomials, orthogonally additive
polynomials, and order-continuous bidual phrases.  The only duplicate was the
earlier orthogonally additive scalar partial packet created by this lane,
which this packet upgrades.  A bounded web search on 2026-06-29 for exact and
near-exact formulations of the scalar upgrade surfaced the source paper but no
later exact answer.

This settles the extracted scalar question.  It does not settle the
vector-valued $G^*$-valued theorem without the source paper's
order-continuity assumptions, and it does not address arbitrary Banach-space
dual isomorphisms that are not lattice isomorphisms.

\begin{thebibliography}{3}

\bibitem{BoydMiranda2025}
C. Boyd and V. C. C. Miranda,
\emph{Lattice isomorphic Banach lattices of polynomials},
arXiv:2509.12417.

\bibitem{BotelhoGarcia2022}
G. Botelho and L. A. Garcia,
\emph{Order continuity of Arens extensions of regular multilinear operators},
arXiv:2201.09116.

\bibitem{BuskesRoberts2019}
G. Buskes and S. Roberts,
\emph{Polynomial valuations on vector lattices},
arXiv:1910.14624.

\end{thebibliography}

\end{document}
